% !PRINTED 423
\weberpage{423}

%<<CONT
\weq{2}{Y^5 + 5aY^2 + 5bY + c = 0,}
worin $a, b, c$ homogene Functionen 3\textsuperscript{ten}, 4\textsuperscript{ten}, 5\textsuperscript{ten} Grades von $\alpha, \beta$ bedeuten, deren Co\"efficienten Ikosa\"ederinvarianten sind.

Die Function $c$ k\"onnen wir leicht aus der Formel bestimmen
\[ c = -\Pi(\tilde\omega)\,\Pi(\alpha + \beta t). \]

Es ist n\"amlich nach (15) des vorigen Paragraphen f\"ur ein unbestimmtes $\lambda$
\[ \lambda^5 - 10f\lambda^3 + 45f^2\lambda - T = \Pi(\lambda - t), \]
also, wenn man $\lambda = -\alpha : \beta$ setzt
\[ \Pi(\alpha + \beta t) = \alpha^5 - 10f\alpha^3\beta^2 + 45f^2\alpha\beta^4 + T\beta^5, \]
und ferner ergiebt sich ohne Weiteres durch Vergleichung eines Gliedes [\S.\,111, (6), (8)]
\[ \Pi(\tilde\omega) = -H^2, \]
also
\weq{3}{c = H^2(\alpha^5 - 10f\alpha^3\beta^2 + 45f^2\alpha\beta^4 + T\beta^5).}

Die Co\"efficienten $a, b$ berechnet man wohl am einfachsten mit H\"ulfe der \gesperrt{Newton}'schen Formeln (Bd.\,I, \S.\,42):
\[
\begin{aligned}
-15a &= S(\alpha\tilde\omega + \beta t\tilde\omega)^3\\
-20b &= S(\alpha\tilde\omega + \beta t\tilde\omega)^4,
\end{aligned}
\]
indem man die Formeln anwendet, die sich nach kurzer Rechnung durch Vergleichung je eines Gliedes der rechten und linken Seite ergeben:
\[
\begin{array}{ll}
S(\tilde\omega^3) = -3.5.8f^2, & S(t\tilde\omega^3) = -5T,\\
S(t^2\tilde\omega^3) = -5.72f^3, & S(t^3\tilde\omega^3) = -3.5fT,\\
S(\tilde\omega^4) = 4.5fH, & S(t^2\tilde\omega^4) = -5.12f^2H,\\
S(t^3\tilde\omega^4) = -5HT, & S(t^4\tilde\omega^4) = -4.5.27f^3H.
\end{array}
\]

So findet sich
\weq{4}{a = 8f^2\alpha^3 + T\alpha^2\beta + 72f^3\alpha\beta^2 + fT\beta^3}
\weq{5}{b = -fH\alpha^4 + 18f^2H\alpha^2\beta^2 + HT\alpha\beta^3 + 27f^3H\beta^4.}

Um daraus die Resolvente der Ikosa\"edergleichung zu erhalten, m\"ussen wir statt $\tilde\omega$ und $t$ die Functionen $u, v$ [\S.\,111, (12)] einf\"uhren.

Setzen wir, indem wir mit $\lambda, \mu$ irgend zwei unbestimmte Gr\"ossen bezeichnen
\weq{6}{\alpha = \frac{\lambda f}{H},\quad \beta = \frac{\mu f^3}{HT},}
%CONT>>
